%%%%%%%%%%%%
%
% $Autor: Wings $
% $Datum: 2019-03-05 08:03:15Z $
% $Pfad: TemplateSensor $
% $Version: 4250 $
% !TeX spellcheck = en_GB/de_DE
% !TeX encoding = utf8
% !TeX root = filename 
% !TeX TXS-program:bibliography = txs:///biber
%
%%%%%%%%%%%%

% Structure
\chapter{Sensor/Actor}

%Introduction
%\Mynote{cite books, applications, board}

\section{Einleitung: Weg- und Positionssensoren}

In der industriellen Fertigungssteuerung nimmt die Erfassung von Zuständen eine zentrale Rolle ein. Gemäß der Fachliteratur wandelt ein Sensor (lat. sensus, der Sinn) „eine physikalische Größe (z. B. Kraft oder Temperatur) mit Hilfe eines physikalischen Effekts in ein elektrisches Signal um“ [Elektrotechnik und Elektronik für Maschinenbauer, S. 433]. Speziell im Maschinenbau wird bei der Steuerung zwischen zwei grundlegenden Messverfahren unterschieden: Der „Positionserfassung durch Schalter“ und der „Positionserfassung durch Wegbeobachtung“ [Elektrotechnik und Elektronik für Maschinenbauer, S. 435].

Bevor der Fokus auf kontaktbehaftete Schalter gelegt wird, sind folgende berührungslose Sensoren (Sensorschalter) als moderner Industriestandard zu nennen:

\begin{itemize}
	\item Induktive Sensoren: Erfassen metallische Elemente „ohne Kontakt zwischen der Bewegung und dem Sensor“ [Elektrotechnik und Elektronik für Maschinenbauer, S. 436].
	\item Kapazitive Sensoren: Nutzen die Beeinflussung eines elektrischen Feldes zur Kapazitätsänderung [Elektrotechnik und Elektronik für Maschinenbauer, S. 438].
	\item Optische Sensoren: Arbeiten mit Infrarot-Lichtstrahlen zur Objekterkennung (z. B. Lichttaster oder Lichtschranken) [Elektrotechnik und Elektronik für Maschinenbauer, S. 440].
	\item Ultraschall-Sensoren: Basieren auf der Messung der Laufzeit von Schallimpulsen [Elektrotechnik und Elektronik für Maschinenbauer, S. 443].
\end{itemize}

\section{Specific Sensor/Actor}

cite board

\section{Specification}

\begin{itemize}
  \item Cite the data sheet
  \item Extract te information from the data sheet
  \item Circuit Diagram
\end{itemize}

\section{Library}

\subsection{Description}

\subsection{Installation}

\begin{itemize}
    \item data types
    \item data structure
    \item data quantity
    \item data quality
\end{itemize}

\subsection{Functions}

\section{Example}

In this section, a simple example is described in detail, which demonstrates how the library supports the sensor/actor.


\subsection{Manual}

\subsection{Inside the Example}

\subsection{Code}

\subsection{Files}



\section{Calibration}

cite method

\section{Simple Code}


\section{Simple Application}



\section{Tests}

\subsection{Simple Function Test}

\subsection{Test all Functions}



\section{Further Readings}

